\documentclass[11pt,reqno]{amsart}
\usepackage{amsmath,amssymb,graphicx,amscd}
\usepackage{mathrsfs}
\usepackage{enumerate}
% pour les S des spheres
\usepackage{mathrsfs}
\usepackage{dsfont}
\usepackage{stmaryrd}
\usepackage{wrapfig}
\usepackage{pstricks}

% ----------------------------------------------------------------
\vfuzz2pt % Don't report over-full v-boxes if over-edge is small
\hfuzz2pt % Don't report over-full h-boxes if over-edge is small

\newcommand{\prob}{{\mathbb{P}}}

\renewcommand{\S}{\mathscr{S}}  % complex unit sphere

\newcommand{\CC}{{\mathbb C}}
\newcommand{\HH}{{\mathbb H}}
\newcommand{\RR}{{\mathbb R}}
\newcommand{\QQ}{{\mathbb{Q}}}
\newcommand{\NN}{{\mathbb{N}}}
\newcommand{\ZZ}{{\mathbb{Z}}}
\renewcommand{\P}{{\mathbb{P}}}
\newcommand{\sgn}{\operatorname{{sgn}}}
 \DeclareMathOperator{\Prob}{\mathbb{P}}   %probability
 \DeclareMathOperator{\E}{\mathbb{E}}      %expectation
 \DeclareMathOperator{\cov}{cov}             %logarithmic integral
\newcommand{\1}{\mathds{1}}

 \newcommand{\I}{1\!\!1}                   %Indicator function


 \newcommand{\rd}{{\mathrm{d}}}            %differential
 \newcommand{\ii}{{\mathrm{i}}}
  \newcommand{\Id}{{\mathrm{Id}}}
\newcommand{\law}{\overset{\mbox{\rm \scriptsize law}}{=}}
\newcommand{\convlaw}{\overset{\mbox{\rm \scriptsize law}}{\longrightarrow}}


\renewcommand{\Re}{{\mathfrak{Re}}}
\renewcommand{\Im}{{\mathfrak{Im}}}

% THEOREMS -------------------------------------------------------

\newtheorem{thm}{Theorem}[section]
\newtheorem{cor}[thm]{Corollary}
\newtheorem{lem}[thm]{Lemma}
\newtheorem{prop}[thm]{Proposition}
\theoremstyle{definition}
\newtheorem{defn}[thm]{Definition}
\theoremstyle{remark}
\newtheorem*{rem}{Remark}



\numberwithin{equation}{section}


\begin{document}

~\vspace{-1cm}



\begin{center}
{\bf Probability, homework 8, due November 11.}
\end{center}
\vspace{0.5cm}
\small




\noindent{\bf Exercise 1.}
Let $(X_1,X_2)$ be a Gaussian vector with mean $(m_1,m_2)$ and non-degenerate covariance matrix $(C_{ij})_{1\leq i,j\leq 2}$.   
Prove that
$$
\E[X_1\mid X_2]=m_1+\frac{C_{12}}{C_{22}}(X_2-m_2).
$$






\noindent{\bf Exercise 2.}
Let $X$ be a random variable such that $\P(X>t)=\exp(-t)$ for any $t\geq 0$. Let $Y=\min(X,s)$, where $s>0$ is fixed.
Prove that, almost surely,
$$
\E[X\mid Y]=Y\1_{Y<s}+(1+s)\1_{Y=s}.
$$

\noindent{\bf Exercise 3}.  Let $\mu$ and $\nu$ be two probability measures such that $\mu\ll \nu$ and $\nu\ll\mu$ (usually abbreviated $\mu\sim\nu$). Let $X=\frac{\rd \mu}{\rd\nu}$.

\begin{enumerate}[(i)]
\item Prove that $\nu(X=0)=0$.
\item Prove that  $ \frac{1}{X}=\frac{\rd \nu}{\rd\mu}$  almost surely (for $\mu$ or $\nu$).
\end{enumerate}




 \noindent{\bf Exercise 4}.   Let $(X_n)_{n\geq 0}$ be defined on  $(\Omega, \mathcal{A}, \mathbb{P})$. Assume this sequence converges in probability (under $\mathbb{P}$) to
$X$. Let $\mathbb{Q}$ be another probability measure on $(\Omega, \mathcal{A})$ assumed to be absolutely continuous w.r.t. $\mathbb{P}$. Prove that
$X_n\to X$ in probability under $\mathbb{Q}$.\\


\noindent{\bf Exercise 5}. Let $Y$ be an integrable random variable on $(\Omega,\mathcal{A},\mathbb{P})$ and $\mathcal{G}$ a sub $\sigma$-field of $\mathcal{A}$. Suppose that $\mathcal{H}\subset\mathcal{G}$ is a sub $\sigma$-field of $\mathcal{G}$. Show that
$\mathbb{E}\left(\mathbb{E}(Y\mid \mathcal{G})\mid\mathcal{H}\right)=\mathbb{E}\left(Y\mid \mathcal{H}\right)$ (almost surely).\\

\noindent{\bf Exercise 6.}
Let $X_1,\dots,X_n$ be i.i.d. integrable random variables, and $S=\sum_{i=1}^n X_i$. Calculate $\E[S\mid X_1]$ and $\E[X_1\mid S]$.\\

\vspace{0.2cm}

\noindent{\bf Exercise 7.}
For fixed $a,b>0$, let $(X,Y)$ be a $\mathbb{N}\times\mathbb{R}_+$-valued random variable such that
$$
\mathbb{P}(X=n,Y\leq t)=b\int_0^t\frac{(ay)^n}{n!}e^{-(a+b)y}\rd y.
$$
For $h:\mathbb{R}_+\to\mathbb{R}$ continuous and bounded, calculate $\E[h(Y)\mid X]$.  Calculate $\E[\frac{Y}{X+1}]$. 
Calculate $\mathbb{P}(X=n\mid Y)$. Calculate $\E[X\mid Y]$.


\vspace{0.2cm}


\noindent{\bf Exercise 8}.  Let $X$ and $Y$ be random variables on $(\Omega,\mathcal{F},\mathbb{P})$, and $\mathcal{G}, \mathcal{H}$ sub $\sigma$-fields of $\mathcal{F}$ such that $\sigma(\mathcal{G},\mathcal{H})=\mathcal{F}$. Find counterexamples to the following assertions:
\begin{enumerate}[(i)]
\item If $\E[X\mid Y]=\E[X]$ then $X$ and $Y$ are independent.
\item If $\E[X\mid \mathcal{G}]=\E[X\mid \mathcal{H}]=0$ then $X=0$.
\item If $X$ and $Y$ are independent then so are $\E[X\mid\mathcal{G}]$ and $\E[Y\mid\mathcal{G}]$.
\end{enumerate}

\vspace{0.2cm}


\noindent{\bf Exercise 9}. Let  $(X_n)_{n\geq 1}$ be a sequence of nonnegative random variables on $(\Omega,\mathcal{A},\mathbb{P})$, and $(\mathcal{F}_n)_{n\geq 0}$ a sequence of sub $\sigma$-fields of $\mathcal{F}$. Assume that $\E(X_n\mid \mathcal{F}_n)$ converges to $0$ in probability. 

\begin{enumerate}[(i)]
\item Show that $X_n$ converges to $0$ in probability.
\item Show that the reciprocal is wrong.
\end{enumerate}

\vspace{0.2cm}




\noindent{\bf Exercise 10.}
On the same probability space,  let $X,Y$ be positive random variables such that $\E[X\mid Y]=Y$ and $\E[Y\mid X]=X$ (almost surely).
Prove that $X=Y$ almost surely.


\noindent{\bf Exercise 11}.
For any $d\geq 1$, we admit that there is only one probability measure $\mu$ on $\mathcal{S}_d$, 
(the $(d-1)$-th dimensional sphere embedded in $\RR^d$) that is uniform, in the following sense: for any isometry $A\in {\rm O}(d)$ (the orthogonal group in $\RR^d$), and any continuous function $f:\mathcal{S}_d\to\RR$,
$$
\int_{\mathcal{S}_d} f(x)\rd\mu(x)=\int_{\mathcal{S}_d} f(A x)\rd\mu(x).
$$
Let $X=(X_1,\dots,X_d)$ be a vector of independent centered and reduced Gaussian random variables.

a) Prove that the random variable $U=X/\|X\|_{L^2}$ is uniformly distributed on the sphere.

b) Prove that, as $d\to\infty$, the main part of the globe is concentrated close to the Equator,
i.e. for any $\varepsilon>0$,
$$
\int_{x\in\mathcal{S}_d,|x_1|<\epsilon} \rd\mu(x)\to 1.
$$

\end{document} 